\section{Топология репозиториев}

В результате выполнения вышеперечисленных команд получились следующие топологии связей в репозиториях. Они оформлены в виде графов при помощи следующих инструментов:
\begin{itemize}
  \item \texttt{svn-graph.pl} --- официальный Perl-скрипт для формирования графов в формате \texttt{Graphviz};
  \item \texttt{git-graph.sh} --- самописный Bash-скрипт для формирования графов в формате \texttt{Graphviz}.
\end{itemize}

\vspace*{\fill}
\begin{figure}[H]
  \centering
  \begin{adjustbox}{max width=\textwidth, max height=0.6\textheight, keepaspectratio}
    \includegraphics{tikz/svn_graph.pdf}
  \end{adjustbox}
  \caption{История удалённого репозитория \texttt{Subversion}}
\end{figure}
\vspace*{\fill}

\begin{figure}
  \centering
  \begin{adjustbox}{max width=0.7\textwidth, max height=0.95\textheight, keepaspectratio}
    \includegraphics{tikz/git_graph.pdf}
  \end{adjustbox}
  \caption{История удалённого репозитория \texttt{Git}}
\end{figure}